\documentclass[11pt,a4paper]{article}

\usepackage{iftex}
\RequirePDFTeX
\usepackage{cmap}
\usepackage[T2A]{fontenc}
\usepackage[utf8]{inputenc}
\usepackage[english,russian]{babel}
\usepackage{PTSerif}
\usepackage{amsmath,amssymb,microtype}
\usepackage[a4paper,left=23mm,right=23mm,top=20mm,bottom=21mm]{geometry}
\usepackage{xcolor,enumitem,fancyhdr,hyperref}

\definecolor{accent}{HTML}{187AA5}
\definecolor{ink}{HTML}{302B28}
\definecolor{muted}{HTML}{6D6560}
\definecolor{tint}{HTML}{EDF4F6}
\definecolor{rulecolor}{HTML}{A9CADA}
\hypersetup{unicode=true,colorlinks=true,urlcolor=accent,linkcolor=accent,
  pdftitle={Домашнее задание 6. Манёвры уклонения от столкновения.},
  pdfauthor={Чиняев Владимир Николаевич}}
\urlstyle{same}
\color{ink}
\setlength{\parindent}{0pt}
\setlength{\parskip}{3pt}
\linespread{1.08}
\setlist[enumerate]{leftmargin=*,label=\textcolor{accent}{\arabic*.},
  itemsep=3pt,topsep=4pt,parsep=0pt}
\setlist[itemize]{leftmargin=*,itemsep=3pt,topsep=3pt,parsep=0pt}
\setlength{\emergencystretch}{1em}
\pagestyle{fancy}
\fancyhf{}
\renewcommand{\headrulewidth}{0pt}
\renewcommand{\footrulewidth}{0.4pt}
\renewcommand{\footrule}{\hbox to\headwidth{\color{rulecolor}\leaders\hrule height\footrulewidth\hfill}}
\fancyfoot[L]{\footnotesize\color{muted}Введение в маневрирование КА\quad /\quad ДЗ 6}
\fancyfoot[R]{\footnotesize\color{muted}\thepage}

\newcommand{\taskheading}[2]{%
  \par\vspace{5pt}{\large\bfseries\color{accent}Задание #1. #2}\par\smallskip}
\newcommand{\hint}[1]{%
  \par\nobreak\smallskip\noindent
  \begingroup
  \setlength{\fboxsep}{7pt}%
  \setlength{\fboxrule}{0.4pt}%
  \fcolorbox{rulecolor}{tint}{%
    \begin{minipage}{\dimexpr\linewidth-2\fboxsep-2\fboxrule\relax}
    \small\textbf{\color{accent}Подсказка.} #1\par
    \end{minipage}}%
  \endgroup\par\smallskip}

\begin{document}

{\small\color{accent}ВВЕДЕНИЕ В МАНЕВРИРОВАНИЕ КА}\par
\vspace{3pt}
{\huge Домашнее задание 6}\par
{\large Манёвры уклонения от столкновения}\par
{\small\color{muted}К лекции 9}
\par{\small\bfseries\color{accent}Срок сдачи: 23 ноября 2026 года, 23:59}
\vspace{4pt}
{\color{rulecolor}\hrule}
\vspace{5pt}

Сравните опасность двух сближений и подберите трансверсальный манёвр
для одного из них. Оформите решение обоих заданий в предоставленном Jupyter Notebook.

\taskheading{1}{Какие сближения опаснее?}

\begingroup
\setlength{\fboxsep}{9pt}
\colorbox{tint}{\begin{minipage}{\dimexpr\linewidth-2\fboxsep\relax}
\textbf{Исходные данные}
\par\smallskip
Рассмотрите два случая номинального сближения без манёвра.
В опорную эпоху $t=0$ относительные положения в инерциальных осях равны
\[
 \boldsymbol\rho_A(0)=\begin{pmatrix}70\\0\\1200\end{pmatrix}\text{м},
 \qquad
 \boldsymbol\rho_B(0)=\begin{pmatrix}0\\110\\1800\end{pmatrix}\text{м}.
\]
В обоих случаях относительная скорость постоянна:
\[
 \boldsymbol v_{\mathrm{rel}}=\begin{pmatrix}0\\0\\-20\end{pmatrix}\text{м/с}.
\]
Ошибки положений объектов независимы, имеют гауссово распределение
с нулевыми математическими ожиданиями. Их ковариации в заданных
инерциальных осях постоянны:
\[
 C_{\mathrm{sc}}=\operatorname{diag}(30^2,60^2,90^2)\ \text{м}^2,
 \qquad
 C_{\mathrm{obj}}=\operatorname{diag}(40^2,80^2,120^2)\ \text{м}^2.
\]
Примите суммарный радиус столкновения и допустимую вероятность
\[
 R_c=1\ \text{м},\qquad P_{\text{доп}}=5\cdot10^{-5}.
\]
\end{minipage}}
\endgroup

Для расчёта вероятности в обоих заданиях используйте приближение
\[
 P_c\approx\frac{R_c^2}{2\sqrt{\det C_b}}
 \exp\!\left(-\frac12\boldsymbol b^\top C_b^{-1}\boldsymbol b\right).
\]

\begin{enumerate}
\item Для каждого случая найдите момент наибольшего сближения и минимальное
расстояние. Оси плоскости сближения направьте вдоль $X,Y$;
запишите матрицу $P$ и найдите векторы промаха.
\item Рассчитайте $C_\rho$, $C_b$ и вероятность столкновения для каждого
случая. Представьте таблицу: время до сближения, минимальное расстояние,
$P_c$, выполнение ограничения $P_c\leq P_{\text{доп}}$.
\item На одном рисунке в плоскости сближения постройте границу
$P_c=P_{\text{доп}}$ и отметьте оба вектора промаха.
Используйте одинаковый масштаб по осям. Объясните, совпадает ли порядок
опасности случаев с порядком минимальных расстояний.
\end{enumerate}

\hint{В прямолинейной модели
$\boldsymbol\rho(t)=\boldsymbol\rho(0)+\boldsymbol v_{\mathrm{rel}}t$.
Момент сближения найдите из условия
$\boldsymbol\rho(t)^\top\boldsymbol v_{\mathrm{rel}}=0$.
Для построения границы возьмите логарифм заданной формулы вероятности.}

\newpage
{\small\color{muted}Домашнее задание 6\hfill К лекции 9}\par
\vspace{5pt}

\taskheading{2}{Когда и в какую сторону выполнять манёвр?}

Для случая \textbf{B} подберите один трансверсальный импульс
в ОСК управляемого КА.
Моменты манёвра отсчитываются от той же опорной эпохи $t=0$:
\[
 \Delta\boldsymbol V=(0,\Delta V_t,0)^\top,\quad
 t_{m,1}=t_{\mathrm{ca}}-600\ \text{с},\quad
 t_{m,2}=t_{\mathrm{ca}}-1800\ \text{с}.
\]
Влияние импульса на вектор промаха уже рассчитано:
\[
\begin{aligned}
 \boldsymbol b(\Delta V_t;t_{m,1})
 &=\boldsymbol b_0+\begin{pmatrix}1200\\600\end{pmatrix}\Delta V_t,\\[4pt]
 \boldsymbol b(\Delta V_t;t_{m,2})
 &=\boldsymbol b_0+\begin{pmatrix}3600\\-1200\end{pmatrix}\Delta V_t.
\end{aligned}
\]
Здесь $\boldsymbol b_0$ --- промах случая B без манёвра.
Коэффициенты имеют размерность секунды; дополнительно умножать их
на время упреждения не нужно. Ковариацию $C_b$, радиус $R_c$
и порог вероятности сохраните из задания 1.
Доступный импульс ограничен:
\[
 |\Delta V_t|\leq0{,}006\ \text{м/с}.
\]

\begin{enumerate}
\item Для каждого момента получите зависимость $P_c(\Delta V_t)$.
Найдите обе границы $P_c=P_{\text{доп}}$ и определите наименьший по модулю
импульс, удовлетворяющий ограничению на вероятность.
Рассмотрите оба знака импульса. На этом этапе ограничение
на величину импульса не накладывайте.
\item Проверьте ограничение на величину импульса и выберите момент
выполнения манёвра. Для выбранного решения укажите компоненты импульса,
характеристическую скорость, вектор промаха и вероятность столкновения.
\item Постройте обе зависимости $P_c(\Delta V_t)$ на одном графике
при $-0{,}04\leq\Delta V_t\leq0{,}04$~м/с. Нанесите порог вероятности,
границы доступного импульса и выбранное решение. Объясните,
как время манёвра повлияло на его величину и знак.
\end{enumerate}

\hint{Подставьте линейную зависимость $\boldsymbol b(\Delta V_t)$
в формулу вероятности и возьмите логарифм равенства
$P_c=P_{\text{доп}}$. Получится квадратное уравнение.
Из двух граничных значений выберите ближайшее к нулю,
затем проверьте ограничение $|\Delta V_t|\leq0{,}006$~м/с.}

\vspace{4pt}
\textbf{Работа с ноутбуком и самопроверка}

В \texttt{hw\_6.ipynb} готовы константы, функции для таблиц и графиков,
чекеры \texttt{check\_task1} и \texttt{check\_task2}.
Заполните отмеченные переменные и текстовые ячейки своими решениями.
Использование чекеров и сохранение результатов их вызова не требуются.
Расчёты выполняйте без промежуточного округления:
положения --- в метрах, время --- в секундах, импульсы --- в м/с.

\vspace{5pt}
{\color{rulecolor}\hrule}
\vspace{5pt}
\textbf{Что представить}

Сдайте один ноутбук с полными решениями обоих заданий:
расчётными формулами, таблицами, обоими рисунками, выбранным манёвром
и объяснениями результатов.

Формулы и рассуждения оформите в текстовых ячейках, расчёты и рисунки ---
в кодовых. Перед отправкой перезапустите ядро, выполните ячейки решения по порядку
и сохраните ноутбук с результатами.

Название файла: \texttt{hw\_6\_name\_surname.ipynb}.
Вместо \texttt{name\_surname} укажите свои имя и фамилию латиницей.

\end{document}
